%% LyX 2.4.4 created this file.  For more info, see https://www.lyx.org/.
%% Do not edit unless you really know what you are doing.
\documentclass[english]{article}
\usepackage[T1]{fontenc}
\usepackage[latin9]{inputenc}
\usepackage{units}
\usepackage{enumitem}
\usepackage{amsmath}
\usepackage{amsthm}
\usepackage{graphicx}
\usepackage{geometry}
\geometry{verbose,tmargin=1in,bmargin=1in,lmargin=1in,rmargin=1in}

\makeatletter
%%%%%%%%%%%%%%%%%%%%%%%%%%%%%% Textclass specific LaTeX commands.
\numberwithin{equation}{section}
\numberwithin{figure}{section}
\newlength{\lyxlabelwidth}      % auxiliary length 

%%%%%%%%%%%%%%%%%%%%%%%%%%%%%% User specified LaTeX commands.
\usepackage{fancyhdr}
\usepackage{lastpage}
\pagestyle{fancy}
\usepackage{enumitem}
\usepackage{ifthen}
%sets math fonts to be more readable
\everymath=\expandafter{\the\everymath\displaystyle}
%\DeclareMathSizes{10}{10}{10}{10}
\usepackage{multicol}
% Added by lyx2lyx
\setlength{\parskip}{\smallskipamount}
\setlength{\parindent}{0pt}

\makeatother

\usepackage{babel}
\begin{document}
\renewcommand{\author}{Dr.\,James Ashe}

\newcommand{\classNum}{232}

\newcommand{\classSec}{C}

\newcommand{\nameType}{}

\newcommand{\examType}{Final Exam Review}

\renewcommand{\date}{Final Exam will be be given on January $16^{\text{th}}$ 10am-12pm, 2012} %%\currenttime, \today

%%First page's header%%%%%%%%%%%%%%%%%%%%%%
\fancypagestyle{firststyle} %%header settings for first pagr
{
	\fancyhead[L]{MTH \classNum{}(\classSec{}) \author{}\\
	\textbf{\examType{}} \date}
	\fancyhead[C]{\textbf{\nameType}}
	\setlength{\headsep}{14pt}
}
%%%%%%%%%%%%%%%%%%%%%%%%%%%%%%%%%%%%%%%%%%

%%Next page's header%%%%%%%%%%%%%%%%%%%%%%%
\setlist{nolistsep}
\lhead{\classNum{}(\classSec{})} 
\chead{\textbf{\nameType}} 
\cfoot{\thepage{}~of~\pageref{LastPage}} 
\setlength{\headsep}{5pt} 
%%%%%%%%%%%%%%%%%%%%%%%%%%%%%%%%%%%%%%%%%%%

%%Various settings; see the preamble for other settings%%%%%
%%\setenumerate[0]{leftmargin=12pt,itemindent=0pt,labelwidth=0pt}
\renewcommand{\labelenumi}{\textbf{\arabic{enumi}.}}
\renewcommand{\labelenumii}{\textbf{\alph{enumii}.}}
\thispagestyle{firststyle}
%%%%%%%%%%%%%%%%%%%%%%%%%%%%%%%%%%%%%%%%%%

\subsubsection*{Riemann Sums}

\emph{Complete the following steps for $f(x)=24-2x$, interval $[0,6]$,
and $n=3$.}
\begin{enumerate}
\item Sketch the graph of the function on the given interval and illustrate
the \emph{right Riemann sums }using $n$ subintervals.
\item Calculate the \emph{right Riemann sums }for the given interval and
number of subintervals, $n$.
\end{enumerate}

\subsubsection*{Understanding Integration}

\emph{Graph an integrand and use geometry to evaluate the integral.}
\begin{enumerate}[resume]
\item ${\displaystyle \int_{1}^{3}2x\,dx}$ (use $\frac{1}{2}bh=area\,of\,a\,triangle$)
\end{enumerate}
\emph{The figure shows the area of regions bounded by the graph of
$f(x)$. Evaluate the following integrals.}

\begin{multicols}{2}
\begin{enumerate}[resume]
\item ${\displaystyle \int_{0}^{8}f}(x)\,dx$
\item ${\displaystyle \int_{0}^{5}f(x)\,dx}-{\displaystyle \int_{8}^{0}f(x)\,dx}$\columnbreak
\end{enumerate}
\includegraphics[width=3cm]{images/Exam_1_area_under_graph.svg}\end{multicols}

\subsubsection*{Evaluating indefinite/definite integrals.}

\begin{multicols}{2}
\begin{enumerate}[resume]
\item [\textbf{6.}]\setcounter{enumi}{6}$f(x)=9x^{8}$

\begin{enumerate}
\item ${\displaystyle \int9x^{8}\,dx}$
\item ${\displaystyle \int_{0}^{1}9x^{8}\,dx}$
\end{enumerate}
\item $f(x)={\displaystyle x^{2}-8x+1}$

\begin{enumerate}
\item $\int(x^{2}-8x+1)\,dx$
\end{enumerate}
\item $f(x)=\cos3x$

\begin{enumerate}
\item ${\displaystyle \int\cos3x\,dx}$
\item ${\displaystyle \int_{-\nicefrac{\pi}{3}}^{\nicefrac{\pi}{3}}\cos3x\,dx}$
\end{enumerate}
\item $f(x)=\frac{x-\sqrt{x}}{x^{3}}$

\begin{enumerate}
\item ${\displaystyle \int\frac{x-\sqrt{x}}{x^{3}}}\,dx$\columnbreak
\end{enumerate}
\item $f(x)=x^{2}(3x^{3}+1)^{4}$

\begin{enumerate}
\item ${\displaystyle \int x^{2}(3x^{3}+1)^{4}\,dx}$
\item ${\displaystyle \int_{1}^{2}x^{2}(3x^{3}+1)^{4}\,dx}$
\end{enumerate}
\item $v(\theta)={\displaystyle \cos\theta\sin^{2}\theta}$

\begin{enumerate}
\item ${\displaystyle \int\cos\theta\sin^{2}\theta\,d\theta}$
\end{enumerate}
\item ${\displaystyle g(t)=\frac{t}{\sqrt{4-t^{2}}}}$

\begin{enumerate}
\item ${\displaystyle \int\frac{t}{\sqrt{4-t^{2}}}\,dt}$
\item ${\displaystyle \int_{2}^{4}\frac{t}{\sqrt{4-t^{2}}}\,dt}$
\end{enumerate}
\end{enumerate}
\end{multicols}

\subsubsection*{Average value of a function over an interval.}
\begin{enumerate}[resume]
\item [\textbf{13.}]\setcounter{enumi}{13}A hiking trail has an elevation
given by $h(x)=40x^{3}-300x^{2}+200x+4000$ where $h$ measures feet
above sea level and $x$ represents miles along the trail with $0\leq x\leq5$.
What is the average elevation along the trail?
\end{enumerate}

\subsubsection*{Initial value problems}

\emph{Solve the initial position problem .}
\begin{enumerate}[resume]
\item $\frac{dS}{dt}=(t+9)^{\nicefrac{1}{2}}$ with $S(0)=20$. Find $S(t)$.
\end{enumerate}
\emph{Solve the problem.}
\begin{enumerate}[resume]
\item A projectile is launched vertically from the ground at $t=0$ where
$t$ is seconds, and its velocity in flight (in $\nicefrac{\mbox{m}}{\mbox{s}}$)
is given by $v(t)=20-10t$.

\begin{enumerate}
\item Find the position at time $t$ for $0\leq t\leq4$.
\item Find the total distance traveled from $t=0$ to $t=4$.
\end{enumerate}
\end{enumerate}

\subsubsection*{Area between graphs}

\emph{Use calculus to find the area of the region described.}
\begin{enumerate}[resume]
\item The region bounded by $y=x$ and $y=x^{2}-4x$.\\
\includegraphics[width=3cm]{images/x_and_x2-4x.png}
\item The region bounded by $y=\frac{1}{2}x^{2}$ and $y=-x^{2}+6$.\\
\includegraphics[width=3cm]{images/1/2x2_and-x2+6.png}
\item The region in the first quadrant bounded by $-x+1$ and $x-1$ from
$0\leq x\leq2$.\\
\includegraphics[width=3cm]{images/1-x_and_x-1.png}
\end{enumerate}

\subsubsection*{Volume of Solids}

\emph{Let $R$ be the region bounded by the following curves. Use
the disk method to find the volume of the solid generated when $R$
is revolved about the $x$-axis.}
\begin{enumerate}[resume]
\item $y=\sqrt{x}$, $x=0$, and $x=9$.
\item $y=\sqrt{\sin4x}$, $y=0$, $0\leq x\leq\frac{\pi}{4}$.
\end{enumerate}

\subsubsection*{Length of the curve}

\emph{Find the length of the curve.}
\begin{enumerate}[resume]
\item $y=2x^{\nicefrac{3}{2}}$ from $x=0$ to $x=\frac{5}{4}$.
\item $y=\left(4-x^{\nicefrac{2}{3}}\right)^{\nicefrac{3}{2}}$ from $x=0$
to $x=8$.
\end{enumerate}
\newpage

\subsubsection*{Derivatives and Integrals of $\ln x$ and $e^{x}$.}

\begin{multicols}{2}
\begin{enumerate}[resume]
\item $\frac{d}{dx}\left(\ln x+1\right)$
\item $\frac{d}{dx}$$3e^{x}\ln x$
\item $\frac{d}{dx}\left(\ln(\cos x)-2x\right)$
\item $\frac{d}{dx}$$e^{-2x+x^{2}}$
\item $\frac{d}{dt}\frac{e^{t}}{e^{t}+1}$
\item ${\displaystyle \int\left(3e^{-3x}+\frac{1}{x}\right)\,dx}$
\end{enumerate}
\end{multicols}

\subsubsection*{Integration by parts}

\begin{multicols}{2}
\begin{enumerate}[resume]
\item [\textbf{29.}]\setcounter{enumi}{29}${\displaystyle \int}xe^{-4x}\,dx$
\item ${\displaystyle \int x\cos4x\,dx}$
\item ${\displaystyle \int\frac{4\ln x}{x^{5}}\,dx}$
\item ${\displaystyle \int_{0}^{\ln1}3xe^{x}\,dx}$ 
\end{enumerate}
\end{multicols}

\subsubsection*{Integrate functions with even/odd powers of $\sin x$ and $\cos x$.}

\begin{multicols}{2}
\begin{enumerate}
\item [\textbf{33.}]\setcounter{enumi}{33}${\displaystyle \int\ln3x\,dx}$
\item ${\displaystyle \int}\sin^{2}3x\,dx$
\item ${\displaystyle \int\cos^{3}x\,dx}$
\item ${\displaystyle \int\sin^{2}x}\cos^{2}x\,dx$
\item ${\displaystyle \int\cos^{3}x}\sin x\,dx$
\end{enumerate}
\end{multicols}

\subsubsection*{Integrate functions use Trigonometric substitution. }

\begin{multicols}{2}
\begin{enumerate}
\item [\textbf{38.}]\setcounter{enumi}{38}${\displaystyle \int\frac{dx}{\sqrt{4-x^{2}}}}$
\item ${\displaystyle \int_{0}^{4}\sqrt{16-x^{2}}\,dx}$
\item ${\displaystyle \int\sqrt{16-x^{2}}}\,dx$
\end{enumerate}
\end{multicols}

\subsubsection*{Method of Partial Fractions.}
\begin{enumerate}[resume]
\item [\textbf{41.}]\setcounter{enumi}{41}${\displaystyle \int\frac{8x+22}{x^{2}+4x-5}\,dx}$
\end{enumerate}
${\displaystyle \int\frac{8x+22}{x^{2}+4x-5}\,dx}={\displaystyle \int\frac{8x+22}{(x-1)(x+5)}\,dx}$
and $\frac{8x+22}{(x-1)(x+5)}=\frac{A}{x-1}+\frac{B}{x+5}$.

$\Rightarrow A(x+5)+B(x-1)=8x+22$

$\Rightarrow Ax+5A+Bx-B=2x+5$, clearly then$Ax+Bx=8x$ and $5A-B=22$. 

Solving the first one for $A$, $A=8-B$ and plugging this into the
second gives, $40-5B-B=22\Rightarrow B=3\Rightarrow A=5$.

Now we have, $\frac{8x+22}{(x-1)(x+5)}=\frac{5}{x-1}+\frac{3}{x+5}$,
and we find ${\displaystyle \int\frac{5}{x-1}+\frac{3}{x+5}\,dx}=5\ln|x-1|+3\ln|x+5|$.
\begin{enumerate}[resume]
\item ${\displaystyle \int\frac{6}{x^{2}+2x-8}\,dx}$
\end{enumerate}

\end{document}
